% !Mode:: “TeX:UTF-8”
% !TeX root = ./main.tex
%%%%%%%%%%%%%%%%%%%%%%%%%%%%%%%%%%%%%%%%%%%%%%%%%%%%%%%%%%%%%%%%%%%%%%%%%%%%%%%%%%%%%%%%
%%%%%%%%%%%%%%%%%%%%%%%%%%%东北林业大学%%理学院%%本科生毕业论文LaTeX模板%%%%%%%%%%%%%%%%%%%%%
%%%%%%%%%%%%%%%%%%%%%%%%%%%%%%%%%%%%%%%%%%%%%%%%%%%%%%%%%%%%%%%%%%%%%%%%%%%%%%%%%%%%%%%%
% Copyright 2022-2025 LiYu
%
% This work may be distributed and/or modified under the
% conditions of the LaTeX Project Public License, either version 1.3
% of this license or (at your option) any later version.
%
%
% The Current Maintainer of this work is Li Yu of collage of science in NEFU.
% If you have any questions about the template, please contact------liy@nefu.edu.cn----
%%%%%%%%%%%%%%%%%%%%%%%%%%%%%%%%%%%%%%%%%%%%%%%%%%%%%%%%%%%%%%%%%%%%%%%%%%%%%%%%%%%%%%%%
%注意编译方法：
%
% Compile with: xelatex -> bibtex -> xelatex -> xelatex
%
%此模板供东北林业大学学生免费试用，版权归理学院李雨老师所有，任何人修改发布请标明以上说明。
%=======================================================================================%
%此模板基于理学院李雨老师的毕业论文Latex模板进行修改, 对原模板进行了封装和简化
%=======================================================================================%
\documentclass{nefuthesis}
%==================================================%
% 导言区, 如果有额外需要使用的包, 请在此处引用
%==================================================%
\usepackage{float}
\usepackage{tabularx}

%==================================================%
% 论文相关信息
%==================================================%
\title{分类任务下证据深度学习的不确定性估计问题}{Uncertainty estimation problem of evidence deep learning in classification tasks}
\author{黄刘诚}{Huang Liucheng}
\studentid{2023113155}
\supervisor{马晓剑}{副教授}{Xiaojian Ma}{Associate Professor}
\subdate{2026年3月} % 论文提交日期
\degree{硕士}{Master} % 本科 or 硕士
\specialty{数学}{Mathematics} % 专业名称
\defensedate{2026年6月} % 论文答辩日期

\begin{document}
% 显示封面
\makecover

% 中文摘要环境, 请在此环境内编写中文摘要
\begin{abstract}[cn][1]{证据深度学习；不确定性估计；分类问题；正则化约束}
这是中文摘要的内容，摘要的内容一般包括研究背景、研究目的、研究方法、研究结果和结论等。摘要应简明扼要，突出重点，
避免使用过于专业的术语和缩略语。中文摘要一般在300字左右。
\end{abstract}

% 英文摘要环境, 请在此环境内编写英文摘要
\begin{abstract}[en][2]{Evidence Deep Learning; Uncertainty Estimation; Classification Tasks; Regularization Constraints}
This is English abstract, the content of the abstract generally includes research background, 
research purpose, research methods, research results and conclusions. The abstract should be concise 
and highlight the key points, avoiding overly professional terms and abbreviations. The English abstract 
is generally about 300 words.
\end{abstract}

% 目录
\content

% 论文正文环境
\begin{thesis}
    % 建议按照以下形式进行正文内容的编写
    % 将所有的章节放在sections文件夹中, 每一章节使用的图片按照pictures->chapter number->image name的形式进行命名
    % 例如：第一章的图片放在pictures->chapter1->image name.png中
    % \input{sections/chapter1.tex}
    % \input{sections/chapter2.tex}
    % \input{sections/chapter3.tex}

\end{thesis} 

% % 附录环境, 请在此环境内编写附录
\begin{Appendix}
    % \input{sections/appendix.tex}
\end{Appendix}

% 结论环境, 请在此环境内编写结论
\begin{conclusion}
    % \input{sections/conclusion.tex}
\end{conclusion}

\references{references.bib} % 参考文献

% 致谢环境, 请在此环境内编写致谢
\begin{acknowledgement}
    % \input{sections/acknowledgement.tex}
\end{acknowledgement}

% % 原创性声明和版权授权书命令
\declaration

\end{document} 